电容式液位测量利用的物理效应是：电容器的电容取决于电极之间的介电质。

\subsection{电容计算}

带介电质的平板电容器电容可由相应公式计算，其中 $A$ 表示电容器面积，$d$ 表示极板间距，$\varepsilon$ 表示介电常数。

本实验使用 Lecher 线的电容。它同样可由相应公式计算，其中 $l$ 表示导线长度，$d$ 表示导线中心之间的距离，$R$ 表示导线半径。

\subsection{LC 振荡回路}

实验中应根据谐振原理，利用 LC 振荡回路进行电容测定。为此，将未知测量电容与已知电感并联。

在假设振荡回路无损耗的情况下，可计算其复阻抗。由该振荡回路阻抗最大这一条件，可以直接推得谐振频率。

\subsection{振荡电路}

为了精确确定谐振频率，通常将谐振器接入振荡电路。电路中的自激振荡器会自动适应决定频率元件的固有频率，例如 LC 振荡回路或石英晶体，并为后续测量和分析提供稳定的周期信号。

\subsection{数字频率测量}

待测频率通过与已知参考时钟比较来确定。具体而言，将测量信号和参考信号转换为数字脉冲，并在规定时间窗口内分别计数。根据记录的脉冲数，可以按相应关系计算待测频率。

\subsection{比较器}

比较器用于将振荡器的模拟信号转换为数字矩形信号。比较器输出电压在 $U_{a,\max}$ 和 $U_{a,\min}$ 之间切换，具体取决于输入电压 $U_1$ 是高于还是低于比较值 $U_2$。

\subsection{Schmitt 触发器}

为了保证稳定的数字信号处理，使用 Schmitt 触发器。与简单比较器不同，它具有两个固定的开关阈值，即上阈值和下阈值。由此产生的滞回特性可以抑制含噪输入信号导致的不必要开关过程，从而实现可靠的脉冲计数以确定频率。

\subsection{系统误差传播}

在许多测量任务中，量 $y$ 不能直接测量，而是由多个影响量通过函数 $y = f(x_1, x_2, \ldots, x_i)$ 得到。由此可以得到系统偏差的传播关系，并估计各个偏差对测量结果的影响。
